\section{模型假设与符号说明}

\subsection{模型假设}

为便于示意，可先给出若干常规假设：
\begin{enumerate}
    \item 配送车辆在规划周期内均可正常投入使用，且中途不考虑临时故障。
    \item 每个客户节点仅服务一次，且对应需求量在规划时刻已知。
    \item 路网通行时间可由历史数据或经验参数估计，并在当前求解周期内视为已知量。
\end{enumerate}

\subsection{主要符号}

\begin{table}[H]
    \centering
    \caption{主要符号说明}
    \label{tab:symbols}
    \begin{tabular}{cc}
        \toprule
        符号 & 含义 \\
        \midrule
        $V$ & 客户节点集合 \\
        $K$ & 配送车辆集合 \\
        $c_{ij}$ & 节点 $i$ 到节点 $j$ 的单位运输成本 \\
        $t_{ij}$ & 节点 $i$ 到节点 $j$ 的行驶时间 \\
        $q_i$ & 节点 $i$ 的配送需求量 \\
        \bottomrule
    \end{tabular}
\end{table}

若正式论文中符号较多，可继续扩展为独立的长表格或附录符号表。表~\ref{tab:symbols} 在模板中主要用于展示表题、表格线和文中交叉引用的默认效果。
